\section{Trabajo relacionado}
\label{sec:related_work}
% \vspace{-1em}
\paragraph{Large Multimodal Models (LMMs)}
Los Large Multimodal Models (LMMs) han evolucionado rápidamente en los últimos años. La investigación reciente se ha centrado en mejorar la capacidad de instruction-following de los LMMs mediante una mejor curación de datos \cite{llavanext2024, li2024llavaov, chen2024sharegpt4v, li2024mvbench, ren2024vista, zhang2024video}, en construir mejores vision encoders capaces de procesar imágenes a resoluciones más altas \cite{llavanext2024, laurenccon2024matters, wang2024qwen2}, y en extender los LMMs a la comprensión intercalada de image \cite{laurenccon2024matters, li2024llava, jiang2024mantis} o video \cite{maaz2023video, lin2023video, zhang2024long}. Para mejorar las capacidades de instruction-following, LLaVA-1.5 \cite{liu2024improved} y LLaVA-NeXT \cite{llavanext2024} aprovechan datos multimodales generados por GPT-4 para un tuning enriquecido con conocimiento. Para optimizar la extracción de características visuales, Idefics2 \cite{laurenccon2024matters} aprovecha la estrategia NaViT \cite{dehghani2023patch} y el vision encoder SigLIP \cite{zhai2023sigmoid} para soportar entradas de imagen con resolución nativa. Qwen2-VL \cite{wang2024qwen2} aplica Multimodal RoPE (M-RoPE) para optimizar la extracción de características para distintos tamaños de entrada.

A pesar de los avances significativos, los LMMs actuales de estado del arte siguen siendo ineficientes para manejar entradas multimodales de contexto largo debido a la complejidad cuadrática de las operaciones de causal self-attention. Para mitigar este problema, métodos como Flamingo \cite{alayrac2022flamingo}, OpenFlamingo \cite{awadalla2023openflamingo}, Idefics \cite{laurenccon2023obelics} y LLaMA 3.2 \cite{dubey2024llama} intercalan capas gated cross-attention con capas self-attention para el modelado vision-text. mPlug-Owl3 \cite{ye2024mplug} desarrolla una hyper-attention layer que integra self-attention con cross-attention para soportar comprensión multimodal. Sin embargo, estos modelos generalmente rinden por debajo de los LMMs basados en Transformer tipo LLaVA. En este estudio identificamos que una posible razón es la falta de actualización de vision tokens en los modelos basados en cross-attention. Nuestro \model aborda esta limitación incorporando capas Mamba adicionales para actualizar vision tokens.


\vspace{-1em}
\paragraph{Efficient LLMs/LMMs}
Para mejorar la eficiencia del razonamiento en LMMs, trabajos recientes \cite{li2023videochat, zhang-etal-2023-video, li2024mvbench, li2023blip, zhang2025llava, shu2024video, shen2024longvu} desarrollan técnicas para reducir la longitud de la secuencia de vision tokens. Video-LLaMA \cite{zhang-etal-2023-video} y VideoChat \cite{li2023videochat, li2024mvbench} usan un Q-Former \cite{li2023blip} para comprimir tokens visuales densos en una secuencia más compacta. LLaVA-Mini \cite{zhang2025llava} reduce aún más cada imagen a un solo token. Video-XL \cite{shu2024video} introduce un Visual Summarization Token para resumir la información visual, mientras que LongVU \cite{shen2024longvu} diseña un mecanismo adaptativo de compresión espaciotemporal para reducir el número de video tokens.

Otra línea de investigación, ejemplificada por la serie de modelos Mamba \cite{gu2023mamba, dao2024transformers}, se centra en desarrollar modelos secuenciales de tiempo lineal para reducir la complejidad. Los modelos basados en Mamba \cite{zuo2024falcon} y los modelos híbridos Mamba-Transformer \cite{lieber2024jamba, glorioso2024zamba} han alcanzado éxitos notables en el ámbito del pure language modelling. Sin embargo, estas arquitecturas eficientes han recibido poca atención en el dominio de la comprensión multimodal. LongLLaVA \cite{wang2024longllava} dio un paso inicial al incorporar el modelo híbrido Jamba \cite{lieber2024jamba} como decoder LMM, aunque su rendimiento en comprensión de video largo sigue siendo subóptimo. Esta limitación probablemente se deba a la ausencia de un híbrido LLM fuerte como modelo base y a la falta de entrenamiento a gran escala con datos multimodales de instruction-following.
